\section{Energy density of the emitted radiation}

In SI units, the instantaneous electromagnetic energy density is
\begin{align}
  w(t,{\bf x})
  =\frac{\epsilon_0}{2}|{\bf E}(t,{\bf x})|^2
  +\frac{1}{2\mu_0}|{\bf B}(t,{\bf x})|^2
  =\frac{\epsilon_0}{2}
  \left(|{\bf E}(t,{\bf x})|^2+c^2|{\bf B}(t,{\bf x})|^2\right).
  \label{eq:standard-energy-density}
\end{align}
With the conventions used in Part~III, the electric and magnetic fields are
the following components of the Faraday tensor:
\begin{align}
  E_i(t,{\bf x})=cF^{i0}(t,{\bf x}),\qquad
  B^i(t,{\bf x})=-\frac12\epsilon_{ijk}F^{jk}(t,{\bf x}).
  \label{eq:fields-from-Faraday}
\end{align}

Part~III defines the Fourier transform with a positive sign in the
exponential,
\begin{align}
  F^{\alpha\beta}(\omega,{\bf x})
  =\frac1{2\pi}\int_{-\infty}^{\infty}dt\,
  e^{i\omega t}F^{\alpha\beta}(t,{\bf x}).
  \label{eq:Faraday-Fourier-convention-observables}
\end{align}
The corresponding inverse transform therefore contains the negative sign,
\begin{align}
  F^{\alpha\beta}(t,{\bf x})
  =\int_{-\infty}^{\infty}d\omega\,
  e^{-i\omega t}F^{\alpha\beta}(\omega,{\bf x}).
  \label{eq:Faraday-inverse-Fourier-observables}
\end{align}
Thus the Fourier transforms of the electric and magnetic fields are obtained
directly from the Fourier-transformed Faraday tensor:
\begin{align}
  E_i(\omega,{\bf x})=cF^{i0}(\omega,{\bf x}),\qquad
  B^i(\omega,{\bf x})=-\frac12\epsilon_{ijk}F^{jk}(\omega,{\bf x}).
  \label{eq:Fourier-fields-from-Faraday}
\end{align}
Substitution of the inverse transforms in
Eq.~\eqref{eq:standard-energy-density} gives the instantaneous energy density
entirely in terms of the available Fourier components:
\begin{align}
  w(t,{\bf x})
  =\frac{\epsilon_0c^2}{2}
  \int_{-\infty}^{\infty}d\omega
  \int_{-\infty}^{\infty}d\omega'\,
  e^{-i(\omega+\omega')t}
  \left[
    \sum_{i=1}^{3}F^{i0}(\omega,{\bf x})F^{i0}(\omega',{\bf x})
    +\sum_{1\leq i<j\leq3}F^{ij}(\omega,{\bf x})F^{ij}(\omega',{\bf x})
  \right].
  \label{eq:instantaneous-energy-density-Fourier-F}
\end{align}

For radiation spectra, the relevant quantity is the energy density integrated
over time.  Parseval's identity for the convention in
Eq.~\eqref{eq:Faraday-Fourier-convention-observables} yields
\begin{align}
  &\int_{-\infty}^{\infty}dt\,w(t,{\bf x})
  =\int_{-\infty}^{\infty}d\omega\,
  \frac{d\mathcal U}{d\omega}(\omega,{\bf x}),\\
  &\frac{d\mathcal U}{d\omega}(\omega,{\bf x})
  =\pi\epsilon_0c^2
  \left[
    \sum_{i=1}^{3}|F^{i0}(\omega,{\bf x})|^2
    +\sum_{1\leq i<j\leq3}|F^{ij}(\omega,{\bf x})|^2
  \right].
  \label{eq:two-sided-spectral-energy-density-F}
\end{align}
This is the two-sided spectrum.  Since the time-domain Faraday tensor is
real, $F^{\alpha\beta}(-\omega,{\bf x})=
F^{\alpha\beta}(\omega,{\bf x})^*$.
\begin{highlightbox}[title={One-sided spectral energy density}]
  Therefore the one-sided spectrum for
  $\omega>0$ is
  \begin{align}
    \frac{d\mathcal U_+}{d\omega}(\omega,{\bf x})
    =2\pi\epsilon_0c^2
    \left[
      \sum_{i=1}^{3}|F^{i0}(\omega,{\bf x})|^2
      +\sum_{1\leq i<j\leq3}|F^{ij}(\omega,{\bf x})|^2
    \right],\qquad \omega>0.
    \label{eq:one-sided-spectral-energy-density-F}
  \end{align}

\end{highlightbox}
